% Pursuing Stacks (cotes 134-2 to 134-8): archive cote 134-8, site page 24 (PDF page 25), typescript page 578.
% typescript. Draft compared: ps7.tex:1110-1169.
% AI-transcribed by Claude Opus 5.5 (claude-opus-5-5, adaptive thinking, effort high), Message Batches API, batch ps4 (msgbatch_012Aoc1NuskATPsyBeKSYK7J), 2026-11-12; not yet checked by eye.
\scanpage{134-8/24}
\mspage{578}

\uline{Remark\hand{s}}\hand{, 1)}One of the (rather few) instances in the duality formulaire
where finiteness conditions are clearly needed,  is the so-called
"biduality theorem",  when taking \uline{RHom}'s with values in a so-called
"dualizing complex", \xout{\unsure{Rxxxxxxx D$^{\mathrm b}$(A}} which here should be an object \xout{in}
\[ R^{\bullet} \quad \text{in} \quad D^{b}(A^{\wedge}_{k}) \]
(for given A in (Cat) and given ring of coefficients\xout{\unsure{,}} k,  for instance
k=$\underline{Z}$  ).  The question arises here, \xout{ax} for any given A and k,  a) whe-%
ther there exists a dualizing complex R$^{\bullet}$  (which,as usual,  will be
unique up to \inserted{dimension shift and} \hand{"}twist\hand{"} by an invertible sheaf of k-modules on A),  and
b) can such dualizing complex be obtained as
\[ R^{\bullet} = p^{!}(k_{e}) \quad , \]
where  p: A $\longrightarrow$ e is the projection to the final object of (Cat) ?  It
should be easy enough,  once the basic duality formalism is written down
for ordered sets,  as contemplated above, \inserted{to show} that when A is a finite or-%
dered set (or only locally finite of finite combinatorial dimension),
then $p^{!}(e_{k})$  is indeed a dualizing complex.  To refer to something
more "en vogue" at present than those poor dualizing complexes,  it is
clear,  from what I heard from Illusie and Mebkhout about the "complexe
d'intersection" for stratified spaces,  that this complex can be descri-%
bed also in the context of locally finite ordered sets (if I got the
stepwise construction of this complex right).  As the construction
here corresponds to stratifications where the strata are by no means
even dimensional,  I am not too sure though if the complex obtained this
way is really relevant  -  it is'nt a topological invariant of the
topological space at any rate,  independently of it's stratification -
a bad point indeed. Too bad !

2)  Here is a rather evident example showing that the asphericity
condition on (13) is not automatic.  Take  I  with a \xout{largest}\inserted[typed]{smallest} element x
(hence  $C_{x} = \hand{|}I\hand{|}$ ),  and J =$\hand{\Delta}_{1}$ \inserted{$= (0 \to 1)$.} \deleted{, consisting of 0 and 1 \inserted{\unsure{= y}} \xout{xxyx}}.  To give
a map  I$\rightarrow$ J  amounts to give  $I_{o} = f^{-1}(0) \subset I$,  which is any \inserted{\unsure{``}}open subset
of I (i.e. such that $a \in I_{o}$ ,  $b \leqslant a$  implies $b \in I_{o}$).  \inserted{If we take} \deleted{The relation}
\deleted{f(x)=y just means that \xout{I$_o$xx$\neq$xI x} I$_o$}.  $I_{o} = \hand{\{}x\hand{\}}$,  then properness of
f is equivalent with  $I_{1} = I \hand{\smallsetminus} I_{o} = I\hand{\smallsetminus}\hand{\{}x\hand{\}}$ being aspheric,  while the
weaker condition contemplated first (the sets (13) non empty) means only
that $I_{1} \neq \emptyset$ .  Now\hand{\unsure{,}} $I_{1}$ \xout{can}\inserted[typed]{may} of course be taken to be any (finite say)
ordered set - the construction made amounting to taking the cone over
the map $\hand{|}I_{1}\hand{|} \longrightarrow e$ .   In this example,  all fibers of $\hand{|}f\hand{|} :\hand{|}I\hand{|}\rightarrow\hand{|}J\hand{|}$
are (canonically) homeomorphic to  $I_{1}$  , except the fiber at \inserted{the "barycenter"} 1, which

\margin{24}
